(1) 以 \((\nu_e,\nu_x)\) 为基，去掉公共能量后 \(H=\frac12[(V-\Delta\cos2\theta)\sigma_z+\Delta\sin2\theta\,\sigma_x]\)。物质势只作用于电子味，公共迹不影响味转化概率。
(2) 两本征值为 \(\pm\Omega/2\)，其中 \(\Omega=\sqrt{(\Delta\cos2\theta-V)^2+\Delta^2\sin^22\theta}\)。定义 \(\sin2\theta_m=\Delta\sin2\theta/\Omega\)、\(\cos2\theta_m=(\Delta\cos2\theta-V)/\Omega\)，则 \(\tan2\theta_m=\Delta\sin2\theta/(\Delta\cos2\theta-V)\)。
(3) \(U(L)=\cos(\Omega L/2)\mathbf1-2iH\sin(\Omega L/2)/\Omega\)。非对角振幅 \(U_{xe}=-i(\Delta\sin2\theta/\Omega)\sin(\Omega L/2)\)，故 \(P_{e\to x}(L)=\sin^22\theta_m\,\sin^2(\Omega L/2)\)，振荡长度 \(L_{\rm osc}=2\pi/\Omega\)。
(4) 当 \(V=\Delta\cos2\theta\) 时 \(\sin^22\theta_m=1\)，为MSW共振；\(V\to0\) 恢复真空振荡。反中微子物质势变号 \(V\to-V\)，若 \(\Delta>0\) 且 \(\cos2\theta>0\)，相同正电子密度下反中微子没有对应共振。
